\documentclass[11pt]{article}
\usepackage[margin=1in]{geometry}
\usepackage{amsmath,amssymb,amsthm}
\usepackage{xurl}
\usepackage{hyperref}
\sloppy
\let\oldus\_ \renewcommand{\_}{\oldus\allowbreak}
\newtheorem{theorem}{Theorem}
\newtheorem{lemma}[theorem]{Lemma}
\theoremstyle{definition}
\newtheorem{definition}[theorem]{Definition}
\newcommand{\Ext}{\mathrel{\prec}}
\title{Conjecture 00000003526: well-founded trees, their rank, and ordinal height}
\author{Jackmeson1}
\date{}
\begin{document}
\maketitle

\section*{Statement}
\emph{Definition: Combinatorial principles of large ordinals: trees and partitions. Conjecture: The
ordinal version of the tree lemma: the vanishing of descending chains of ordinal trees is
well-foundedness, with the rank of well-foundedness the ordinal height.}

\paragraph{Verdict.} PROVED, for the two standard notions of tree named below. The claim is classical.
We do not take it from a library theorem about trees. We write out the definitions of tree, branch,
descending chain, rank and height, and we prove every equivalence from them. All of it is checked in
Lean 4 with Mathlib.

\section{Readings and definitions}
\textbf{Reading A (trees of finite sequences).} We use the trees of descriptive set theory. Wikipedia,
``Tree (descriptive set theory)'', says that a tree on $X$ is a set of finite sequences ``such that
every prefix of a sequence in the collection also belongs to the collection''. It says that ``a tree
is a nonempty subset'' of $X^{<\omega}$ with this closure property, and that it ``shows that the
empty sequence belongs to every tree''. A branch is an infinite sequence ``each of whose finite
prefixes belongs to'' $T$, and ``A tree that has no branches is called wellfounded; a tree with at
least one branch is illfounded.'' An \emph{ordinal tree} is a tree on a set of ordinals. All our
results hold for trees on an arbitrary set $A$, so they hold in particular for $A=\mathrm{Ord}$.

\textbf{Reading B (set-theoretic trees).} Wikipedia, ``Tree (set theory)'', says that ``a tree is a
partially ordered set'' $(T,<)$ such that for each $t$ the set $\{s:s<t\}$ ``is well-ordered by the
relation'' $<$. The order type of that set ``is called the height, rank, or level of'' $t$, and the
height of $T$ ``itself is the least ordinal greater than the height of each element of'' $T$.

\begin{definition}[Lean names in brackets]
Let $T\subseteq A^{<\omega}$.
\begin{itemize}
\item $T$ is a \emph{tree} [\texttt{IsTree}] if $T\neq\emptyset$ and $t\in T$, $m<|t|$ imply
  $t{\restriction}m\in T$.
\item For $f:\mathbb N\to A$, $f{\restriction}n=\langle f(0),\dots,f(n-1)\rangle$ [\texttt{initSeg}].
  The sequence $f$ is a \emph{branch} [\texttt{IsBranch}] if $f{\restriction}n\in T$ for all $n$. The
  tree $T$ is \emph{well-founded} [\texttt{TreeWF}] if it has no branch.
\item \emph{Strict extension} [\texttt{Ext}]: for nodes $s,t\in T$, $t\Ext s$ iff $s$ is a proper
  prefix of $t$. A \emph{descending chain} is a sequence of nodes with $c_{n+1}\Ext c_n$ for all $n$,
  that is, $c_0\subsetneq c_1\subsetneq\cdots$ [\texttt{HasDescChain}]. This is the orientation in which
  ``descending'' means going away from the root.
\item If $\Ext$ is well-founded, the \emph{rank} [\texttt{nodeRank}] is
  $\rho_T(s)=\sup\{\rho_T(t)+1:t\in T,\ t\Ext s\}$, with $\sup\emptyset=0$. This is Mathlib's
  well-founded rank \texttt{Acc.rank}. The \emph{rank of the tree} [\texttt{treeRank}] is
  $\rho(T)=\sup\{\rho_T(s)+1:s\in T\}$.
\item A \emph{ranking} [\texttt{IsRanking}] is a map $\varphi:T\to\mathrm{Ord}$ with
  $t\Ext s\Rightarrow\varphi(t)<\varphi(s)$. The \emph{ordinal height} [\texttt{ordHeight}] is
  $\mathrm{ht}(T)=\min\{\alpha:\exists\varphi\text{ a ranking with }\varphi(s)<\alpha\ \forall s\}$.
  This definition does not mention $\rho$.
\end{itemize}
\end{definition}
For the name ``ordinal height'': Huschenbett, Kartzow, Liu and Lohrey, \emph{Tree-Automatic
Well-Founded Trees} (arXiv:1201.5495), Section 2.2, write ``Let us recall the (ordinal) rank of a
well-founded partial order.'' They define the rank of an element by recursion and then ``the rank
of'' the order ``(also called the ordinal height of'' it) as $\sup\{\mathrm{rank}(p)+1\}$. This is
exactly $\rho(T)$. Our $\mathrm{ht}(T)$ is a second, independent notion of height, and Theorem 3 shows
that the two agree.

\paragraph{Conventions.} Ordinals live in Lean's universe-polymorphic \texttt{Ordinal.\{u\}}. For a
tree on $A$ with $A$ in \texttt{Type u}, ranks lie in \texttt{Ordinal.\{u\}}. For a tree on
\texttt{Ordinal.\{u\}} itself, ranks lie in \texttt{Ordinal.\{u+1\}}, and $\alpha$ appears as
$\mathrm{lift}(\alpha)$. In Lean, \texttt{nodeRank} is $0$ on ill-founded trees and \texttt{ordHeight}
is $\inf\emptyset=0$ when no ranking exists. These junk values never enter a theorem, because every
rank statement assumes well-foundedness. The step from ``no descending chain'' to ``well-founded'' uses
dependent choice, which Lean's \texttt{Classical.choice} provides. As Wikipedia, ``Well-founded
relation'', puts it: ``Equivalently, assuming the axiom of dependent choice, a relation is
well-founded when it contains no infinite descending chains, meaning there is no infinite sequence
x0, x1, x2, ... of elements of X such that xn+1 R xn for every natural number n.''

\section{Reading A}
\begin{theorem}[\texttt{treeWF\_iff\_no\_descChain}, \texttt{no\_descChain\_iff\_wf}]
For a tree $T$, the following are equivalent: $T$ has no branch; $T$ has no descending chain;
$\Ext$ is well-founded on $T$.
\end{theorem}
\begin{proof}
Let $f$ be a branch. Then $c_n=f{\restriction}n$ is a descending chain, because
$f{\restriction}(n+1)=f{\restriction}n\frown f(n)$. Conversely, let $(c_n)$ be a descending chain. By
induction, $c_m\subseteq c_n$ for $m\le n$. Since a proper prefix is strictly shorter,
$|c_n|\ge n$. Put $f(k)=c_{k+1}(k)$. For $i<n$ we have $c_{i+1}\subseteq c_n$, so
$f(i)=c_n(i)$. Hence $f{\restriction}n=c_n{\restriction}n$, which is a prefix of $c_n\in T$. By
prefix closure (\texttt{IsTree.prefix\_mem}) it lies in $T$, so $f$ is a branch. The second
equivalence is the descending-chain criterion
(\texttt{wellFounded\_iff\_isEmpty\_descending\_chain}).
\end{proof}

\begin{theorem}[\texttt{nodeRank\_eq\_iSup\_child}, \texttt{nodeRank\_unique}, \texttt{nodeRank\_le\_of\_ranking}, \texttt{wf\_iff\_exists\_ranking}]
Let $T$ be a tree. $\Ext$ is well-founded if and only if $T$ has a ranking. If $T$ is well-founded,
then:
(a) $\rho_T(s)=\sup\{\rho_T(s\frown a)+1: s\frown a\in T\}$, the supremum over one-step extensions;
(b) $\rho_T$ is the only function satisfying the recursion (a);
(c) $\rho_T(s)\le\varphi(s)$ for every ranking $\varphi$, and $\rho_T$ is itself a ranking.
\end{theorem}
\begin{proof}
Suppose $\varphi$ is a ranking. Then $\Ext$ is a subrelation of the inverse image of $<$ on the
ordinals under $\varphi$, so it is well-founded. Conversely, $\rho_T$ is a ranking because
$t\Ext s$ implies $\rho_T(t)<\rho_T(s)$.

(a) Every one-step extension $s\frown a\in T$ satisfies $s\frown a\Ext s$, which gives $\ge$. For
$\le$, take $t\Ext s$. Then $t=s\frown a\frown w$, so $s\frown a$ is a prefix of $t$ and lies in $T$.
Hence $\rho_T(t)\le\rho_T(s\frown a)$, with equality if $t=s\frown a$.

(b) Let $\varphi$ satisfy the recursion. Well-founded induction along $\Ext$, using (a), shows
$\varphi(s)=\rho_T(s)$.

(c) By well-founded induction: if $\rho_T(t)\le\varphi(t)<\varphi(s)$ for all $t\Ext s$, then
$\rho_T(s)=\sup_{t\Ext s}(\rho_T(t)+1)\le\varphi(s)$.
\end{proof}

\begin{theorem}[\texttt{ordHeight\_eq\_treeRank}, \texttt{treeRank\_eq\_succ\_root}]
If the tree $T$ is well-founded, then $\mathrm{ht}(T)=\rho(T)=\rho_T(\emptyset)+1$.
\end{theorem}
\begin{proof}
$\rho_T$ is a ranking with values below $\rho(T)$, so $\mathrm{ht}(T)\le\rho(T)$. Let $\varphi$ be a
ranking bounded by $\alpha$. By Theorem 2(c), $\rho_T(s)+1\le\varphi(s)+1\le\alpha$ for every $s$, so
$\rho(T)\le\alpha$. Every $s$ extends $\emptyset$, so $\rho_T(s)\le\rho_T(\emptyset)$, and therefore
$\rho(T)=\rho_T(\emptyset)+1$.
\end{proof}
The Lean statement \texttt{tree\_wf\_rank\_height} bundles Theorems 1--3 except the uniqueness of the one-step recursion (Theorem~2(b)), which is proved separately as \texttt{nodeRank\_unique}.

\begin{theorem}[\texttt{nodeRank\_decTree}, \texttt{ordinal\_tree\_realization}]
For every ordinal $\alpha$, let $D_\alpha=\{s:\alpha>s_0>s_1>\cdots>s_k\}$ be the tree on the ordinals
of strictly decreasing sequences below $\alpha$ (\texttt{decTree}). Then $D_\alpha$ is a
well-founded tree, $\rho_{D_\alpha}(s)$ is the last entry of $\alpha\frown s$, so
$\rho_{D_\alpha}(\emptyset)=\alpha$, and $\mathrm{ht}(D_\alpha)=\alpha+1$.
\end{theorem}
\begin{proof}
Let $b(s)$ be the last entry of $\alpha\frown s$. We have $s\frown a\in D_\alpha$ iff $s\in D_\alpha$
and $a<b(s)$. Along a proper extension, $b$ strictly decreases, so $b$ is a ranking and $D_\alpha$ is
well-founded. Moreover $\sup\{a+1:a<b(s)\}=b(s)$, so $b$ satisfies recursion (a), and Theorem 2(b)
gives $\rho=b$. The rest follows from Theorem 3.
\end{proof}
So every ordinal, and in particular every countable ordinal, is the root rank of a well-founded
ordinal tree. Every successor ordinal is the height $\mathrm{ht}=\rho(T)$ of such a tree. Under this
convention the height of a nonempty tree is always a successor ordinal (Theorem 3).

\section{Reading B}
\begin{theorem}[\texttt{setTree\_wf\_rank\_height}]
Let $(P,<)$ be a set-theoretic tree [\texttt{IsSetTree}]. Then $<$ is well-founded, so descending
chains vanish. The well-founded rank of each $t$ [\texttt{setRank}] equals its height, the order type
of $\{s:s<t\}$ [\texttt{ht}]. The height of $P$, the least ordinal above all heights
[\texttt{setTreeHeight}], equals $\sup_t(\mathrm{rank}(t)+1)$.
\end{theorem}
\begin{proof}
Suppose $f(n+1)<f(n)$ for all $n$. Then $(f(n+1))_n$ is a descending sequence in the well-order
$\{s:s<f(0)\}$, which is impossible. For $s<t$, the map $z\mapsto z$ is an order isomorphism from
$\{z:z<s\}$ onto the initial segment of $\{z:z<t\}$ below $s$, so $\mathrm{ht}(s)=\mathrm{typein}(s)$.
By induction,
$\mathrm{rank}(t)=\sup_{s<t}(\mathrm{typein}(s)+1)$. Every typein is below $\mathrm{type}\{z:z<t\}$,
and every smaller ordinal is a typein, so this supremum equals $\mathrm{ht}(t)$. The statement about
the height of $P$ follows.
\end{proof}

\section*{Lean formalization and axiom audit}
The file is \texttt{lean/Conjecture3526/Basic.lean} (383 lines, \texttt{import Mathlib} only). Its main
theorems are \texttt{C3526.tree\_wf\_rank\_height}, \texttt{C3526.ordinal\_tree\_realization} and
\texttt{C3526.setTree\_wf\_rank\_height}. The command \texttt{\#print axioms} reports only
\texttt{propext}, \texttt{Classical.choice} and \texttt{Quot.sound} for each of them. The file contains
no \texttt{sorry} and no \texttt{native\_decide}.

\section*{Sources}
\begin{itemize}
\item Wikipedia, Tree (descriptive set theory), \url{https://en.wikipedia.org/wiki/Tree_(descriptive_set_theory)}.
\item Wikipedia, Tree (set theory), \url{https://en.wikipedia.org/wiki/Tree_(set_theory)}.
\item Wikipedia, Well-founded relation, \url{https://en.wikipedia.org/wiki/Well-founded_relation}.
\item M. Huschenbett, A. Kartzow, J. Liu, M. Lohrey, Tree-Automatic Well-Founded Trees, \url{https://arxiv.org/abs/1201.5495}.
\end{itemize}
\end{document}
